% ============================================================================
% BENCH MEASUREMENTS - the only file that needs editing after the lab session.
%
% Each entry below is a slot. Replace the \bench{...} with the value measured,
% keeping the braces. Anything still holding a \bench{...} renders red in the
% PDF, so the report is finished exactly when nothing here is red any more.
%
% Do not put a value here that was not actually measured.
% ============================================================================

% --- Section 2: wiring confirmed with a meter -------------------------------
% One or two sentences: which pins were checked, what the anchor pins (EEPROM
% VCC to 3V3, VSS to GND) did, and whether the README's mapping held.
\newcommand{\mWiring}{\bench{continuity check result}}

% --- Section 3: the SCK and MOSI capture ------------------------------------
\newcommand{\mSckPeriod}{\bench{SCK period, e.g. 4.02 us}}
\newcommand{\mSckFreq}{\bench{1/period, e.g. 248.8 kHz}}
\newcommand{\mSckError}{\bench{percent error against 250 kHz}}
\newcommand{\mClockIdle}{\bench{idle state of SCK, high or low}}
\newcommand{\mClockCycles}{\bench{clock cycles counted in the capture}}

% Bit sequence read off the screen, MSB first, and the value it reconstructs.
% Read this from the waveform, not from the source.
\newcommand{\mDecodedBits}{\bench{eight bits in order}}
\newcommand{\mDecodedValue}{\bench{value they give}}
\newcommand{\mSampleEdge}{\bench{which edge the device samples on}}

% --- Section 4: the EEPROM ---------------------------------------------------
% How the address width was confirmed on the bench, including what the status
% register showed after a deliberate short write if that was tried.
\newcommand{\mAddressWidth}{\bench{how the two-byte address was confirmed}}

\newcommand{\mStatusBefore}{\bench{eeprom\_status\_before}}
\newcommand{\mStatusAfter}{\bench{eeprom\_status\_after}}
\newcommand{\mWriteWait}{\bench{eeprom\_write\_wait\_ms}}
\newcommand{\mValueWritten}{\bench{value written}}
\newcommand{\mValueRead}{\bench{value read back}}
\newcommand{\mVerifyResult}{\bench{passed or failed}}
\newcommand{\mPulseCount}{\bench{clock pulses while CS is low}}

% What was written, what the reset was, what came back afterwards.
\newcommand{\mPersistence}{\bench{persistence test result}}

% --- Section 5: the fault ----------------------------------------------------
% The four table columns. Keep each to a phrase; the narrative goes below.
\newcommand{\mFaultExpected}{\bench{expected}}
\newcommand{\mFaultObserved}{\bench{observed}}
\newcommand{\mFaultCause}{\bench{register and field}}
\newcommand{\mFaultCorrection}{\bench{correction}}

% The diagnosis in prose: the first meaningful difference between the two
% captures, and why the corrected configuration now satisfies the device
% electrically. This is where the marks are.
\newcommand{\mFaultNarrative}{\bench{diagnosis: first difference, why the
correction works electrically}}

% --- Section 6: the state machine --------------------------------------------
% What PC13 did in each mode, with ee_use_fsm at 1 and at 0.
\newcommand{\mHeartbeat}{\bench{heartbeat under the FSM versus the blocking path}}

% One complete run through the state machine.
\newcommand{\mFsmResult}{\bench{one write, read and verify result via the FSM}}
